\chapter{Background and Related Work}
\label{ch:background}

This chapter establishes the technical foundations on which the thesis builds.
Section~\ref{sec:riscv} introduces the RISC-V ISA and the rationale for vector
acceleration. Section~\ref{sec:faults} reviews the physical origins and the
abstract modelling of hardware faults. Section~\ref{sec:redundancy} surveys the
classical hardware redundancy techniques, focusing on the duplication- and
triplication-based schemes used in this work. Section~\ref{sec:ijtag-bg}
introduces the IEEE~1149.1 and IEEE~1687 standards for test and in-field
instrument access. Section~\ref{sec:related} positions the thesis with respect
to related work.

\section{RISC-V and vector acceleration}
\label{sec:riscv}

\subsection{The RISC-V instruction set architecture}

RISC-V~\cite{riscv-isa} is an open, royalty-free instruction set architecture organised as a
small mandatory base integer ISA (\texttt{RV32I} / \texttt{RV64I}) together with
a set of optional standard extensions --- integer multiplication and division
(\texttt{M}), atomics (\texttt{A}), single- and double-precision floating point
(\texttt{F}/\texttt{D}), compressed instructions (\texttt{C}), and the vector
extension (\texttt{V})~\cite{riscv-vector}, among others. Crucially, the ISA explicitly reserves
encoding space for \emph{non-standard, custom extensions}. This is what makes
RISC-V so attractive for domain-specific acceleration: a designer can add
application-specific instructions and the hardware to execute them, while
remaining within a coherent and well-documented architectural framework.

The openness of the ISA has a second consequence that is directly relevant to
this thesis. Because the architecture and many high-quality implementations are
open, researchers can modify the microarchitecture freely --- inserting
redundancy, instrumentation, or monitoring logic at any level of the hierarchy
--- something that is simply impossible with the closed, encrypted RTL of most
commercial cores.

\subsection{Data-level parallelism and SIMD accelerators}

A large fraction of the workloads that run on embedded and edge devices ---
digital filtering, fast transforms, linear algebra, and the convolutions at the
heart of neural networks --- are characterised by \emph{data-level parallelism}:
the same operation is applied independently to many data elements. The natural
hardware answer is the Single-Instruction-Multiple-Data (SIMD) paradigm, in
which one instruction drives a wide datapath that processes many elements
simultaneously.

A SIMD vector accelerator amortises the cost of instruction fetch, decode, and
control over a large amount of arithmetic, which is why such accelerators
deliver excellent energy efficiency on suitable workloads. The flip side, and
the reason this thesis focuses on them, is that the resulting wide arithmetic
datapath --- arrays of adders, multipliers, shifters, and comparators ---
typically dominates the area and the switching activity of the core. From a
reliability standpoint the accelerator is therefore both the \emph{largest}
target a fault can strike and the \emph{busiest}, so faults are activated into
errors with high probability.

\subsection{The Klessydra cores}

Klessydra~\cite{klessydra,klessydra-dsp} is a family of RISC-V cores designed around interleaved hardware
multithreading and configurable SIMD acceleration, intended for research and for
embedded use. The cores implement the RISC-V base ISA together with custom
vector instructions that are executed by a dedicated coprocessor, the DSP unit
studied here. Key configuration parameters --- the SIMD width, the number of
parallel accelerator lanes, the scratchpad-memory size, the data width, and the
thread-pool size --- are exposed as VHDL generics and package constants, so that
a single source base can be elaborated into a wide range of design points. This
parameterisation is exploited heavily in the present work: the same accelerator
source is synthesised under many different redundancy and replication settings.

\section{Faults, errors, and failures}
\label{sec:faults}

\subsection{Terminology}

It is useful to distinguish three levels following the classical dependability~\cite{avizienis}
taxonomy. A \emph{fault} is the physical or logical defect --- a flipped storage
bit, a stuck-at node, a momentary voltage glitch. An \emph{error} is the
manifestation of a fault in the system state --- an incorrect value held in a
register or propagating along a wire. A \emph{failure} occurs when an error
propagates to the system output and the delivered service deviates from the
specification. Fault tolerance is the art of breaking the chain from fault to
failure: either by \emph{detecting} an error before it becomes a failure, or by
\emph{masking} it so that it never reaches the output.

\subsection{Physical origins}

Several physical mechanisms produce faults in modern integrated circuits:

\begin{itemize}
  \item \textbf{Radiation-induced single-event effects (SEEs).} An energetic
        particle --- a cosmic-ray secondary neutron at terrestrial altitude, or
        a heavy ion in space --- deposits charge as it traverses the silicon. If
        the collected charge exceeds a node's critical charge $Q_{crit}$, the
        node's logic value is disturbed. When this disturbance flips a storage
        element it is a \emph{single-event upset} (SEU)\footnote{See~\cite{baumann} for a survey of radiation-induced soft errors in advanced technologies.}; when it produces a
        transient pulse in combinational logic that may be latched downstream it
        is a \emph{single-event transient} (SET).
  \item \textbf{Supply and timing noise.} Voltage droop, ground bounce, and
        crosstalk can momentarily corrupt logic values or cause a path to miss
        its timing, producing transient errors that are strongly workload- and
        environment-dependent.
  \item \textbf{Ageing and wear-out.} Bias-temperature instability,
        hot-carrier injection, and electromigration slowly shift device
        parameters over the service life, eventually turning marginal paths into
        permanent timing failures or open/short defects.
\end{itemize}

As $Q_{crit}$ falls with each technology generation, the susceptibility of an
individual node to transient upset tends to rise, even as the smaller node also
presents a smaller cross-section to an incident particle. The net effect is that
soft errors are a first-order reliability concern across the whole range of
technology nodes considered in this thesis.

\subsection{Fault models}

For design and verification purposes the physical mechanisms above are abstracted
into \emph{fault models}. The \emph{stuck-at} model, in which a node is
permanently fixed at \texttt{0} or \texttt{1}, underpins most manufacturing
test. The \emph{single-event-upset} model, a transient bit-flip of a storage
element, is the workhorse for soft-error analysis. \emph{Transition} and
\emph{path-delay} models capture timing-related and ageing faults. The
redundancy mechanisms developed in this thesis are agnostic to the precise model:
DMR detects any discrepancy between two copies, and TMR masks any single faulty
copy, regardless of whether the underlying cause was a particle strike, a glitch,
or a wear-out defect that has manifested in only one replica.

\section{Hardware redundancy techniques}
\label{sec:redundancy}

Fault tolerance is achieved by adding \emph{redundancy} --- in hardware, in
time~\cite{nicolaidis}, or in information. This thesis concentrates on \emph{spatial (hardware)
redundancy}, the dominant choice for protecting high-throughput datapaths where
re-execution in time would be too costly.

\subsection{Information redundancy}

Information redundancy adds check bits to data so that errors can be detected or
corrected. Parity detects single-bit errors; single-error-correcting,
double-error-detecting (SEC-DED) Hamming codes~\cite{ecc-hamming} are the standard protection for
memories and register files. Information redundancy is extremely efficient for
\emph{storage}, but it is awkward for arbitrary \emph{combinational arithmetic}:
protecting a multiplier with an arithmetic residue code, for instance, requires
a bespoke checker and does not generalise across functional units. For this
reason information codes are complementary to, rather than a replacement for, the
modular redundancy studied here.

\subsection{Dual Modular Redundancy (DMR)}

Dual Modular Redundancy duplicates a module and compares the two outputs
bit-by-bit. If the copies agree, the result is forwarded and no error is
flagged; if they disagree, an error flag is raised. DMR therefore provides
\emph{error detection} but not correction: it cannot, on its own, decide which of
the two copies is correct. Its appeal is its low cost --- roughly twice the logic
of the protected module plus a thin comparison layer --- and its generality: the
same comparator works for any module, regardless of its internal function. In a
system context the DMR error flag can trigger a recovery action such as
re-execution, rollback, or a transition to a safe state.

\subsection{Triple Modular Redundancy (TMR)}

Triple Modular Redundancy~\cite{vonneumann,koren} triplicates a module and feeds the three outputs into a
bitwise majority voter, whose output equals the value produced by at least two of
the three copies. As long as no more than one copy is faulty at any given bit,
the voter \emph{masks} the fault completely: the error never reaches the output
and no recovery action is required. The logic function of a single-bit
three-input majority voter is
\begin{equation}
  v = (a \wedge b) \vee (b \wedge c) \vee (a \wedge c),
  \label{eq:majority}
\end{equation}
where $a$, $b$, $c$ are the three replica outputs. TMR is the canonical choice
for systems that must continue operating correctly through a fault without
interruption, at the cost of roughly three times the logic of the protected
module plus the voter. The voter itself is a single point of failure and in the
most demanding applications is itself triplicated; in this work a single voter
per output bit is used, which is the common practice at the datapath level.

\subsection{Granularity and the cost of protection}

A key design decision --- and a central theme of this thesis --- is the
\emph{granularity} at which redundancy is applied. Redundancy can wrap a single
functional unit, a group of units, or an entire accelerator. Finer granularity
allows the designer to spend the area budget where it matters most (for example,
protecting only the multiplier, which is both large and arithmetically
intricate), while coarser granularity is simpler to deploy and protects more of
the design. Because the functional units differ enormously in size, the cost of
protection depends strongly on \emph{which} unit is hardened, not only on
\emph{whether} DMR or TMR is used --- a point that the experimental results in
Chapter~\ref{ch:results} make quantitatively explicit.

\subsection{Replication for throughput versus replication for reliability}

There is an important and often overlooked interaction between redundancy and
\emph{performance} replication. A SIMD accelerator may already instantiate
multiple parallel lanes to increase throughput. When such replication is already
present, the incremental cost of \emph{also} using those lanes for reliability
--- comparing or voting their results --- can be far smaller than building
redundancy from scratch, because the expensive datapath copies already exist.
This thesis evaluates both regimes: building redundancy around a single-lane
accelerator (Scenario~A), and adding voting on top of an already-replicated
multi-lane accelerator (Scenario~B).

\section{Test and in-field monitoring: JTAG and IJTAG}
\label{sec:ijtag-bg}

\subsection{IEEE 1149.1 (JTAG)}

The IEEE~1149.1 standard~\cite{jtag}, universally known as JTAG, defines a Test Access Port
(TAP) and boundary-scan architecture originally created to test the
interconnections of boards after assembly. A four- or five-wire serial interface
(\texttt{TCK}, \texttt{TMS}, \texttt{TDI}, \texttt{TDO}, and optionally
\texttt{TRST}) drives a finite-state machine that shifts data through scan
registers. Over time JTAG became the de-facto standard not just for boundary
scan but for debug access and on-chip instrument control.

\subsection{IEEE 1687 (IJTAG)}

As chips accumulated more and more embedded instruments --- built-in self-test
controllers, sensors, trimming and configuration registers --- the flat,
fixed JTAG data-register structure became unwieldy. IEEE~1687, or
\emph{Internal JTAG} (IJTAG), standardises~\cite{ijtag} a \emph{reconfigurable, hierarchical}
network for accessing these embedded instruments. Its central element is the
\emph{Segment Insertion Bit} (SIB), a small controllable element that can switch
a segment of the scan path --- and the instruments on it --- into or out of the
active chain. By nesting SIBs, the access network can be reconfigured on the fly
so that only the instruments of interest are included in the scan path, keeping
access times short even when hundreds of instruments are present.

IJTAG is accompanied by two description languages: the \emph{Instrument
Connectivity Language} (ICL), which describes the structure of the network and
the instruments, and the \emph{Procedural Description Language} (PDL), which
describes how to operate them. Commercial tools --- in this work, Siemens~\cite{tessent}
Tessent --- consume ICL to automatically synthesise the network hardware~\cite{ijtag-network} and the
retargetable access patterns.

\subsection{From manufacturing test to in-field monitoring}

Although born for manufacturing test, the IJTAG network is equally a vehicle for
\emph{in-field monitoring}. Embedded sensors --- thermal monitors, voltage and
activity sensors, ageing monitors --- can be attached to the IJTAG network as
instruments, and their readings retrieved at run time through the same standard
interface. This makes IJTAG a natural complement to the structural redundancy of
DMR and TMR: redundancy keeps the design \emph{correct} in the presence of
faults, while IJTAG-attached sensors keep it \emph{observable}, allowing a system
to track temperature, activity, and degradation over its service life. Combining
the two on the same accelerator is one of the goals of this thesis.

\section{Related work}
\label{sec:related}

Hardware redundancy for processors and accelerators is a mature field. TMR has
been used for decades in aerospace and space-grade designs, and lockstep
dual-core arrangements (a coarse form of DMR with comparison) are standard in
automotive safety microcontrollers~\cite{lockstep}. Reliable-systems design is treated comprehensively in~\cite{siewiorek,koren}. At finer granularity, selective hardening
--- protecting only the most vulnerable or most critical sub-blocks --- has been
explored extensively as a way to trade reliability against area and power, often
guided by fault-injection campaigns that rank sub-blocks by their contribution
to the overall failure rate~\cite{mukherjee,riscv-hardening}.

For RISC-V specifically, a growing body of work applies these techniques to open
cores: lockstep and TMR variants of small cores, ECC-protected register files
and memories~\cite{tmr-fpga}, and selectively hardened execution units. Vector and SIMD
accelerators, however, have received comparatively less attention despite their
outsized contribution to area and activity, and systematic, \emph{cross-node}
cost characterisations at the granularity of individual functional units are
scarce.

On the monitoring side, IEEE~1687 has been studied primarily as a test-access
mechanism, with more recent work using it to retarget access to in-field sensors
and BIST engines. The automated insertion of sensors into an HDL hierarchy,
coupled with automatic generation of the IJTAG network, is comparatively
under-explored, particularly in combination with structural redundancy on the
same design.

This thesis contributes to both strands by treating a single, parameterised
RISC-V vector accelerator as a common substrate on which (i) DMR and TMR are
applied at controllable granularity and characterised consistently across three
technology nodes, and (ii) an automated IJTAG sensor-insertion flow provides
in-field observability --- bringing reliability and monitoring together in one
reusable methodology.
